% ==============================================================================
% SLIDES - SEMANA 02: CONTROLE DE VERSÃO AVANÇADO COM GIT (GITFLOW, TRUNK-BASED DEVELOPMENT)
% ==============================================================================

% 1. CAPA
\senaicover
  {Controle de Versão Avançado com Git (GitFlow, Trunk-Based Development)}
  {Gerenciamento de código fonte e estratégias de ramificação colaborativa.}
  {Unidade Curricular: Integração e Entrega Contínua - DevOps}
  {Prof. André Souza}

% 2. AGENDA
\begin{senaiframe}{Visão Geral \& Agenda}{Estrutura da Aula}
  \begin{columns}[t]
    \begin{column}{0.48\textwidth}
      \begin{tcolorbox}[senaibluecard, title={1. Conceitos Chave}]
        \begin{itemize}
          \item Fundamentos teóricos e definições.
          \item Relevância para a Indústria 4.0.
        \end{itemize}
      \end{tcolorbox}
      \vspace{4pt}
      \begin{tcolorbox}[senaiambercard, title={3. Aplicação Prática}]
        \begin{itemize}
          \item Estudo de caso na Célula Smart N1.
          \item Desenvolvimento de rotinas técnicas.
        \end{itemize}
      \end{tcolorbox}
    \end{column}
    \begin{column}{0.48\textwidth}
      \begin{tcolorbox}[senairedcard, title={2. Detalhamento Técnico}]
        \begin{itemize}
          \item Ferramentas e protocolos utilizados.
          \item Boas práticas de implementação.
        \end{itemize}
      \end{tcolorbox}
      \vspace{4pt}
      \begin{tcolorbox}[senaigreencard, title={4. Exercícios \& Avaliação}]
        \begin{itemize}
          \item Desafios de fixação do conhecimento.
          \item Integração com a plataforma FactoryHub.
        \end{itemize}
      \end{tcolorbox}
    \end{column}
  \end{columns}
\end{senaiframe}

% 3. CONTEÚDO TEÓRICO
\begin{senaiframe}{Teoria • Módulo 1}{Controle de Versão Avançado com Git (GitFlow, Trunk-Based Development)}
  \begin{columns}[t]
    \begin{column}{0.48\textwidth}
      \begin{tcolorbox}[senaibluecard, title={Fundamentos Teóricos}]
        \begin{itemize}
          \item Commits atômicos, convenção de mensagens (Conventional Commits), branches (feature, release, hotfix), Pull Requests e resolução de conflitos.
        \end{itemize}
      \end{tcolorbox}
    \end{column}
    \begin{column}{0.48\textwidth}
      \begin{tcolorbox}[senaigreencard, title={Prática \& Laboratório}]
        \begin{itemize}
          \item Exercício prático de simulação de conflito de código e merge usando GitFlow em equipe.
        \end{itemize}
      \end{tcolorbox}
    \end{column}
  \end{columns}
\end{senaiframe}
